\input{../tex-inputs/latex-korrekturansicht-vorspann}

\section[Arthur Schnitzler an Theodor Herzl, 4. 8. 1893]{L03904 Arthur Schnitzler an Theodor Herzl, 4. 8. 1893}
\nopagebreak\mylabel{L03904v}
\rehead{ }\normalsize\beginnumbering\briefempfaengerindex{Herzl, Theodor@\textsc{Herzl, Theodor}!zzzSchnitzler, Arthur@\emph{von Arthur Schnitzler}!1893-08-041@{4. 8. 1893}|(be}
\toendnotes[C]{\smallbreak\pagebreak[2]}
\correspDesc{Versand  durch Arthur Schnitzler am 4. 8. 1893 in Wien
\newline{}Erhalt  durch Theodor Herzl im Zeitraum [4. 8. 1893
                  – 7. 8. 1893?] in Wien}\toendnotes[C]{\smallbreak}
\Standort{Jerusalem, Central Zionist Archives, H1:1924-9.}
\physDesc{Brief, 1 Blatt, 3 Seiten, 877 Zeichen (Briefpapier mit Trauerrand)
\newline{}Handschrift: schwarze Tinte, deutsche Kurrent
\newline{}Ordnung: mit Bleistift von unbekannter Hand innerhalb das Konvoluts paginiert:
                                    »31« }\toendnotes[C]{\smallbreak}
\pstart
           \raggedleft{}{\pb}\uline{\textcolor{pink}{Wien}\oindex{Wien@\textbf{Wien}, \emph{Verwaltungsgebiet}|pw}{}\ledrightnote{\textcolor{pink}{Wien}}{ }4. 8. 93}\pend
           
\pstart{}Verehrteſter Freund,\pend\vspace{0.5em}
\pstart
           haben Sie meinen \label{K_L03904-1v}\edtext{Brief}{\lemma{\textnormal{\emph{Brief}}}\Cendnote{\textnormal{Arthur Schnitzler an Theodor Herzl, 1. 7. 1893.}}}\label{K_L03904-1} in die \textcolor{pink}{Schweiz}\oindex{Schweiz@\textbf{Schweiz}|pw}{}\ledrightnote{\textcolor{pink}{Schweiz}} beko{\geminationm}en? – Waren Sie in
                  \textcolor{pink}{Oeſterreich}\oindex{Österreich@\textbf{Österreich}|pw}{}\ledrightnote{\textcolor{pink}{Österreich}}? – Waren Sie in \textcolor{pink}{Wien}\oindex{Wien@\textbf{Wien}, \emph{Verwaltungsgebiet}|pw}{}\ledrightnote{\textcolor{pink}{Wien}}? Hab ich Sie verſäumt, während ich in \textcolor{pink}{Iſchl}\oindex{Bad Ischl@\textbf{Bad Ischl}|pw}{}\ledrightnote{\textcolor{pink}{Bad Ischl}} war? – Ich denke, es iſt irgend eine Nachricht verloren
               gegangen. Nun läßt mir das genaue {\pb}\label{K_L03904-2v}\edtext{Studium der \textcolor{green}{N. Fr. Pr.}\pwindex{Neue Freie Presse@\emph{Neue Freie Presse}|pw}{}\ledrightnote{\textcolor{green}{Neue Freie Presse}} keinen Zweifel mehr übrig, dß Sie längſt wieder in
                  \textcolor{pink}{Frankreich}\oindex{Frankreich@\textbf{Frankreich}|pw}{}\ledrightnote{\textcolor{pink}{Frankreich}}}{\lemma{\textnormal{\emph{Studium … Frankreich}}}\Cendnote{\textnormal{Er dürfte sich auf das Feuilleton vom
                  gleichen Tag beziehen: \textcolor{blue}{Theodor Herzl}\pwindex{Herzl, Theodor 2.\,5.\,1860 Budapest – 3.\,7.\,1904 Edlach@\textsc{Herzl, Theodor} (2.\,5.\,1860 Budapest – 3.\,7.\,1904 Edlach), \emph{Schriftsteller, Journalist}|pwk}: \emph{\textcolor{green}{Wahlbilder aus Frankreich. I. Die Brücke von Neuville}\pwindex{Herzl, Theodor 2.\,5.\,1860 Budapest – 3.\,7.\,1904 Edlach@\textsc{Herzl, Theodor} (2.\,5.\,1860 Budapest – 3.\,7.\,1904 Edlach), \emph{Schriftsteller, Journalist}!Wahlbilder aus Frankreich. I. Die Brücke von Neuville@\strich\emph{Wahlbilder aus Frankreich. I. Die Brücke von Neuville}|pwk}}.
                     In: \emph{\textcolor{green}{Neue Freie Presse}\pwindex{Neue Freie Presse@\emph{Neue Freie Presse}|pwk}}, Nr. 10.398,
                        4. 8. 1893, Morgenblatt, S. 1–3. }}}\label{K_L03904-2} sind. –\pend
           
\pstart
           — Ich gehe vielleicht \label{K_L03904-3v}\edtext{am 20. Auguſt wieder auf 8–10 Tage weg}{\lemma{\textnormal{\emph{am … weg}}}\Cendnote{\textnormal{Zwischen 22. 8. 1893 und 31. 8. 1893 war \textcolor{blue}{Schnitzler} in
                     \textcolor{pink}{Südtirol}\oindex{Südtirol@\textbf{Südtirol}, \emph{Verwaltungsgebiet}|pwk}, \textcolor{pink}{Kärnten}\oindex{Kärnten@\textbf{Kärnten}, \emph{Land}|pwk}, der \textcolor{pink}{Steiermark}\oindex{Steiermark@\textbf{Steiermark}, \emph{Land}|pwk} und \textcolor{pink}{Niederösterreich}\oindex{Niederösterreich@\textbf{Niederösterreich}, \emph{Land}|pwk}. }}}\label{K_L03904-3}.– Für alle Fälle
               hoff ich bald wieder mit 2 Worten zu erfahren, daß Sie u all d\textcolor{gray}{ie}
               Ihren wohl ſind.–\pend
           
\pstart
           Von meiner Sti{\geminationm}ung will ich lieber gar nicht reden; –
               der Strohhalm, mit dem ich mich an die Lebensfreude kla{\geminationm}ere, {\pb}iſt augenblicklich das \textsc{Bicycle}! – Haben Sie ſich auf Ihrer Reiſe nicht mit dem Drama eingelassen? –
               Sich nicht ausgeruht, indem Sie ein neues Lustspiel ſchrieben? Refrain: Für alle
               Fälle hoffe ich bald wieder {\dots}{ }\textsc{etc.}{ }\textsc{etc.}–\pend
           
\pstart
           Herzlichen Gruß!{\\[\baselineskip]}Ganz der Ihre{\\[\baselineskip]}\spacefill\mbox{ArthurSchnitzler}\pend
           \leftskip=0em{}\selectlanguage{ngerman}\endnumbering\briefempfaengerindex{Herzl, Theodor@\textsc{Herzl, Theodor}!zzzSchnitzler, Arthur@\emph{von Arthur Schnitzler}!1893-08-041@{4. 8. 1893}|)be}\mylabel{L03904h}  \newcommand{\dateiname}{L03904}\newcommand{\titel}{Arthur Schnitzler an Theodor Herzl, 4. 8. 1893}\newcommand{\editorInnen}{Herausgegeben von Jahnke, SelmaMüller, Martin Anton}\input{../tex-inputs/latex-korrekturansicht-abspann}
